% Generated draft: requires mathematical and source review.
\documentclass[11pt]{article}
\usepackage[margin=25mm]{geometry}
\usepackage{fontspec}
\setmainfont{Latin Modern Roman}
\usepackage{amsmath,amssymb,mathtools,unicode-math}
\setmathfont{Latin Modern Math}
\usepackage{enumitem,needspace,xurl,hyperref}
\setlength{\parindent}{0pt}
\setlength{\parskip}{5pt}
\setlength{\emergencystretch}{3em}
\newcommand{\N}{\mathbb{N}}
\newcommand{\Z}{\mathbb{Z}}
\newcommand{\Q}{\mathbb{Q}}
\newcommand{\R}{\mathbb{R}}
\newcommand{\C}{\mathbb{C}}
\newcommand{\F}{\mathbb{F}}
\newcommand{\A}{\mathbb{A}}
\newcommand{\PP}{\mathbb{P}}
\newcommand{\E}{\mathbb{E}}
\newcommand{\eps}{\varepsilon}
\newcommand{\dd}{\,\mathrm{d}}
\begin{document}
\section*{UnsolvedMath ID: 701595}
\textbf{Source problem code:} catalogue entry
\subsection{Definitions and mathematical statement}
Let \(k\) be a field of characteristic \(p>2\), and let \(K/k\) be a finite field extension of degree \(n=[K:k]\geq 1\). Regard \(K\) as an \(n\)-dimensional vector space over \(k\), and fix an ordered \(k\)-basis \(\mathcal B\). Write \(\operatorname{GL}_k(K)\) for the group of invertible \(k\)-linear maps on \(K\); here the source notation \(\operatorname{Aut}_k(K)\) denotes linear automorphisms, not necessarily field automorphisms. For \(t\in K^{\times}=K\setminus\{0\}\), define \(m_t(x)=tx\), and put
\[
T_{\mathcal B}(K/k)=\{[m_t]_{\mathcal B}:t\in K^{\times}\}\leq\operatorname{GL}(n,k).
\]
This is the subgroup denoted by \(T(K/k)\) and called a ``nonsplit maximal torus'' in the supplied source. No separability hypothesis is supplied; the explicit multiplication subgroup is the definition used here, without asserting an algebraic-group torus interpretation in every case.

For distinct indices \(i,j\in\{1,\ldots,n\}\), let \(E_{ij}\) be the matrix with entry \(1\) in position \((i,j)\) and zero elsewhere. The elementary transvections relative to \(\mathcal B\) are
\[
u_{ij}(a)=I_n+aE_{ij},\qquad a\in k\setminus\{0\},\quad i\ne j.
\]
A subgroup \(H\leq\operatorname{GL}(n,k)\) is \emph{rich in transvections} if \(u_{ij}(a)\in H\) for every such \(i,j,a\). When \(n=1\), the set of distinct index pairs is empty; the source supplies no restriction \(n\geq2\).

The source asks whether every subgroup \(H\leq\operatorname{GL}(n,k)\) containing \(T_{\mathcal B}(K/k)\) and a ``one-dimensional transformation'' is rich in transvections. The supplied record does not define ``one-dimensional transformation,'' so this hypothesis remains unresolved. If the intended meaning is an invertible transformation \(g\) with \(\operatorname{rank}_k(g-I_n)=1\), then the corresponding provisional formulation is: for every \(k,K,\mathcal B\) as above and every \(H\leq\operatorname{GL}(n,k)\), is the implication
\[
\left(T_{\mathcal B}(K/k)\leq H\ \text{and}\ \exists g\in H:\operatorname{rank}_k(g-I_n)=1\right)
\Longrightarrow
\left(\forall i\ne j\ \forall a\in k\setminus\{0\},\ I_n+aE_{ij}\in H\right)
\]
valid? This rank interpretation is conditional and requires source confirmation; it is not a supplied definition.
\subsection{Short English statement}
For a finite extension of a field of odd characteristic, consider the subgroup of invertible linear maps given by multiplication by nonzero extension elements. Must every linear group containing this subgroup and a ``one-dimensional transformation'' contain every elementary transvection? The meaning of ``one-dimensional transformation'' is missing; the displayed rank formulation is provisional.
\subsection{Sources}
\begin{itemize}
\item Local UnsolvedMath registry, permanent ID \(701595\), exact source problem code \texttt{catalogue entry}, title ``Kourovka Notebook Problem 19.47.'' The supplied registry status is \texttt{open}; this is a source label, not an independently verified literature assessment.
\item Supplied source attribution: \emph{The Kourovka Notebook (Unsolved Problems in Group Theory)}, 19th issue (2018), problem 19.47; the background also identifies the 21st issue edition (2026), arXiv:1401.0300v48. These references have not been independently consulted. The collection description names E. I. Khukhro and V. D. Mazurov as editors.
\item The supplied proposer field is malformed and reads \texttt{V. A. Ko\textbackslash u\{\textbackslash i}. The background repeats this fragment and places \texttt{baev\}} after the heading ``Discussion and literature.'' The intended proposer spelling cannot be recovered reliably from this record. The supplied proposed year is 2018.
\item The record was created and updated on \texttt{2026-10-16T14:00:00Z}, with upstream revision \texttt{d371ee917ac9647fbb06cd6fc529deec6f3dda82}. These are registry timestamps, not a supplied literature-verification date. No \texttt{literature\_checked\_at} value or \texttt{literature\_sources} links are supplied.
\end{itemize}
\end{document}
